\begin{corollary}[The dilatation component from the pair kernels]
\label{cor:c10-P7}
Let a family of gauge-invariant cocycles stable under $\Lambda$ satisfy the hypotheses of Definition \ref{def:c10-hyp}, with the solution $\hbar F_1^{(1)}(\Delta)$ whose quadratic part is \eqref{eq:c9-F1Delta} at the computed values $\theta=\theta_\lambda=\theta_V=-\tfrac13$ of Theorem \ref{thm:c9-theta2}, the translation-invariant solution of Theorem \ref{thm:c9-quadratic} (iii) minus $\tfrac13$ units times $S^{(1)}(\Delta)$. If on its cohomology the composite term of $O$ minus the bracket of $O$ with the translation-invariant solution minus $4$ units times the field-counting functional induces $-\tfrac13$ units times $\{S^{(1)}(\Delta),-\}$, then the spectrum of $D_1$ on the family is $4$ times that of the map of Theorem \ref{thm:c10-brst}: twice the one-loop anomalous dimensions. The condition holds on the part with the outputs $A|_{W_{\mathbb R}^\perp}$ alone of the scalar bilinears (Theorem \ref{thm:c9-D1}) and of the planar scalar sector (Proposition \ref{prop:c10-scalarsector}), on the cohomology of the scalar bilinears when their parts with the outputs $(A|_V,A|_V)$ and $(c,A^*|_V)$ lie in the image of $\Lambda^{(0)}$ (Corollary \ref{cor:c9-ident} (iii)), and on the families of Theorem \ref{thm:c11-families}, where it is the pair of identities (b), (c) of Definition \ref{def:c11-P7} (Theorem \ref{thm:c11-P7} (i)).
\end{corollary}

\begin{proof}
By Proposition \ref{prop:c10-spec} (ii) the term $-\tfrac13$ units times $S^{(1)}(\Delta)$ of the solution contributes $+\tfrac13$ units times $\{S^{(1)}(\Delta),-\}$, which cancels the operator of the condition. $S^{(1)}(\Delta)$ is annihilated by $Q$ and by $\{I,-\}$ (Lemma \ref{lem:c9-SDelta} (iii)), so this solution solves \eqref{eq:c9-arity1} at every order in the fields; Theorem \ref{thm:c9-theta2} identifies it as the solution compatible with the arity-two equations on two legs.
\end{proof}
